\section{How to Run}

From the \texttt{smart-home-alarm-system} directory generated by the CI,
which contains \texttt{pom.xml}:

\begin{itemize}
  \item Build: \texttt{mvn -f pom.xml compile}
  \item Test: \texttt{mvn -f pom.xml test}
  \item CLI: \texttt{mvn -f pom.xml exec:java "-Dexec.mainClass=pcd.shas.Main"}
  \item Demo: \texttt{mvn -f pom.xml exec:java "-Dexec.mainClass=pcd.shas.DemoMain"}
\end{itemize}

The CLI entry point is implemented in
\texttt{src/main/java/pcd/shas/Main.java}; the demo entry point is in
\texttt{src/main/java/pcd/shas/DemoMain.java}.

The interactive CLI accepts commands such as \texttt{arm full PIN}, \texttt{arm partial PIN PERIMETER GROUND\_FLOOR},
\texttt{pin PIN}, \texttt{front door}, \texttt{ground floor}, \texttt{living room}, \texttt{bedroom},
\texttt{status}, and \texttt{quit}.
The default configuration uses an exit delay and an entry delay of five seconds.

The demo scenario uses the same actor graph with shorter delays.
It exercises full arming, entry delay, alarm activation, PIN-based disarming, partial arming, and ignored inactive-zone
sensors.
